\section*{物理：}

加速度不变，相同位移通过时间之比值：

弹簧串联 $k=$\sbl \quad 弹簧并联 $k=$\sbl

弹簧弹性势能 $W=$\mbl

绳子中间系活结，弹力方向为\sbl 两边绳大小\sbl

渡河，水速大于船速是，距离最短的方向：

平抛，速度偏向角是位移偏向角的\sbl

向心力：$F=$\sbl $=$\sbl $=$\sbl

火车轮内\sbl 外\sbl ，轨道内\sbl 外\sbl

开普勒第一定律：

开普勒第二定律：

开普勒第三定律：

万有引力公式：

第一宇宙速度意义：

第二宇宙速度意义：

动量 $P=$\sbl $=$\sbl （和动能 $E_k$ 关系）$E_k=$\sbl

动物（$v_1$）撞静物（$v_2$）：

\column{$v_1^{'} =$\mbl $v_1$ \sep $v_2^{'} =$\mbl $v_1$
}
机械振动：

弹簧/简谐运动力：$F=$\sbl $=$\sbl ，可推得 $ma=$\sbl

\column{$x=$\xlbl \sep $v=$\xlbl \sep $a=$\xlbl}

$E_p=$\xlbl $=$\xlbl

$E_k=$\xlbl $=$\xlbl

$E_{p+k}=$\xlbl

波：\xlbl \hspace{2em} 符号意义：

单摆：要求摆动幅度\mbl

回复力：$F=$\lbl \hspace{2em} 周期：$T=$\lbl

波跨介质传播，周期\sbl 变，频率\sbl 变

波从 A 介质（入射角$\theta_1$）射入 B 介质（入射角$\theta_2$）

\mbl $=$\mbl $=$相对折射率 $n$

折射率 $n=$\mbl $=$\mbl

衍射：

衍射条件：

全反射：

全反射条件：

双缝干涉本质与公式：

双份干涉亮条纹：\lbl \quad 暗条纹\clb  \quad 间距\clb

干涉：

\column{加强点：\lbl \sep 减弱点：\lbl \sep 强度：\lbl}

\subsection*{电学}

电场力：$F=$\sbl ，$k=$\sbl ，$\varepsilon=$\sbl ，静电力做功 $W=$\sbl

电场强度 $E=$\sbl $=$\sbl

电势 $\phi=$\sbl ，$U_{AB}=$\sbl

电容 $C=$\sbl \quad 平行板电容器电容 $C=$\sbl

并联电容\sbl \quad 串联电容\sbl

电流 $I=$\sbl $=$\sbl $=$\sbl \quad 电动势 $E=$\sbl $=$\sbl

电阻 $R=$\sbl $=$\sbl

电功 $W=$\sbl $=$\sbl ，纯电阻电路 $Q=$\sbl

“改表”：有 G 表和 R 电阻，画接线方式、电压大小、扩大倍数

\column{电压表 \sep 电流表}

\subsection*{X 手定则：}

\column{ \sep 安培定则 \sep 洛伦兹力 \sep 电磁感应}

求什么

哪只手

代表什么

计算公式 

\subsection*{电与磁}

磁感应强度：\sbl ，单位\sbl $=$\sbl $=$\sbl

磁通量：\sbl ，单位\sbl $=$\sbl

$B=$\sbl （与力）$=$\sbl （与 $k$）$=$\sbl （与磁通量）$F_{\text{洛}}=$\sbl

$\Phi=$\sbl

带电粒子在磁场中：

$F_{\text{向}}=$\sbl \quad 半径 $R=$\sbl $=$\sbl

周期 $T=$\sbl \quad 频率 $f=$\sbl $=$\sbl

角速度 $\omega=$\sbl $=$\sbl

感应电动势瞬时值：\sbl \quad 发电机线圈转动角速度：\sbl \quad 感应电动势最大值：\sbl

变压器：$U_1 : U_2 =$\sbl \quad $I_1 : I_2 =$\sbl

电磁感应感应电动势 $E=$\sbl $=$\sbl

理想气体体积方程\clb

\section*{振动与波}

单摆适用条件：$\theta <$\sbl

单摆周期公式：$T=$\mbl

受迫振动：驱动力频率 $f_{\text{驱}} =$\sbl 时发生共振，振幅\sbl

阻尼振动：振幅逐渐\sbl （因摩擦/阻力做功）

减幅振动受力分析：一动力方向相同（助力），一相反（阻力），合力\sbl

\subsection*{光的干涉、衍射与偏振}

单缝衍射特点：中央亮纹最\sbl 最\sbl ，两侧对称分布明暗条纹

双缝干涉相邻明纹间距：$\Delta x =$\mbl

频率 $f \uparrow \rightarrow$ 波长 $\lambda $\sbl $\rightarrow \Delta x $\sbl $\rightarrow$ 条纹变\sbl

红光波长长 $\rightarrow$ 条纹\sbl ；紫光波长短 $\rightarrow$ 条纹\sbl

薄膜干涉：
肥皂膜、油膜 $\rightarrow$ \sbl 表面反射光干涉
厚度为 \sbl 或 \sbl 时出现同一颜色条纹
白光照射呈现\sbl 条纹；单色光呈现\sbl 相间条纹

增透膜原理：利用相消干涉减少\sbl 光

劈尖干涉：
条纹凸向薄处说明待测面是\sbl ；凸向厚处说明待测面是\sbl
（亮暗均匀分布）

偏振：\sbl 波特有性质，纵波无偏振
自然光通过偏振片 $\rightarrow$ 成为\sbl 偏振光
马吕斯定律：$I =$\mbl
应用举例：\sbl 、\sbl 、\sbl

\subsection*{光的本性：波粒二象性}

体现波动性的现象：\sbl 、\sbl 、\sbl
体现粒子性的现象：\sbl 、\sbl 、\sbl

光子能量：$E =$\sbl $=$\sbl
光子动量：$p =$\sbl $=$\sbl
康普顿散射证实光子具有\sbl

\subsection*{全反射}

临界角公式：$\sin C =$\sbl （$n_1 > n_2$）
从介质射入空气时：$\sin C =$\sbl

\subsection*{分子动理论与热学}

分子运动速率分布图特点：中间\sbl 两头\sbl
增加温度，分子平均速率\sbl ，但并非所有分子速率都增加

理想气体状态方程：$C =$\sbl （常数）
液体毛细现象：浸润 $\rightarrow$ 液面中间\sbl ；不浸润 $\rightarrow$ 液面中间\sbl

黑体辐射：
黑体吸收所有入射电磁波，不\sbl 也不\sbl
黑体辐射谱仅与\sbl 有关！

光电效应：
爱因斯坦解释 $\rightarrow$ 光子能量 $E =$\sbl
逸出功 $W_0 =$\sbl
最大初动能 $E_{k,\max} =$\sbl
截止电压 $U_c =$\sbl
光电子发射条件：$\nu $\sbl $\nu_0$ （即 $\lambda $\sbl $\lambda_0$）
饱和电流 $\propto$ \sbl
反向电压使电流减小至零 $\rightarrow$ \sbl 电压

光电管实验：
加正向电压/增加光强 $\rightarrow$ \sbl 光电流
加反向电压/减小光强 $\rightarrow$ \sbl 光电流

氢原子能级跃迁：
高 $\rightarrow$ 低：\sbl 光子
低 $\rightarrow$ 高：\sbl 光子 / 受电子撞击
吸收光子能量必须严格等于两能级\sbl
从第 $N$ 能级向下跃迁，发射光子种类：\sbl 种
例如：从 $n=4 \rightarrow$ 基态，可发出\sbl 种不同频率光子
能级图：$n=1,2,3,...$ 能量递增（负值绝对值\sbl ）
电离条件：$E $\sbl $0 \rightarrow$ 自由电子
$n=3,4,5,6 \rightarrow n=2$ 能级的光在\sbl 范围内
频率正比于\sbl ，反比于\sbl

德布罗意波（物质波）：
任何物体都有\sbl 性
统一波长公式：$\lambda =$\sbl （$p$ 为动量）
波长增加，衍射现象\sbl

\subsection*{原子核物理}

核反应方程守恒律：
1. \sbl 守恒
2. \sbl 守恒

衰变类型：
$\alpha$ 衰变 $\rightarrow$ \sbl 粒子（He核），带\sbl 电，电离\sbl ，穿透\sbl
$\beta$ 衰变 $\rightarrow$ \sbl 粒子（电子），带\sbl 电，电离\sbl ，穿透\sbl
$\gamma$ 衰变 $\rightarrow$ \sbl 射线（高能光子），伴随 $\alpha/\beta$ 衰变，电离\sbl ，穿透\sbl

半衰期规律只在\sbl 原子时符合

重核裂变（质量数 $>60$）
轻核聚变

稳定性判断：不是结合能越大越稳定，而是\sbl 越大越稳定
比结合能定义：\sbl / 核子数
比结合能在质量数 \sbl 附近达到最大值

质能方程：$E =$\sbl
基本电荷：$1e =$\sbl C

\subsection*{热力学定律}

理想气体内能：忽略分子势能，内能只与\sbl 有关
单原子分子内能：$U =$\sbl $nRT$
理想气体状态方程：$PV =$\sbl
摄氏温标与热力学温标转换：$K = t + $\sbl （一般取 273）

热力学第一定律：$\Delta U =$\sbl $+$\sbl
符号规定：吸热 $\Rightarrow Q $\sbl ；放热 $\Rightarrow Q $\sbl
气体膨胀 $\rightarrow$ 对外做功 $\rightarrow W $\sbl $0$
气体压缩 $\rightarrow$ 外界对气体做功 $\rightarrow W $\sbl $0$

$p-V$ 图中面积代表：\sbl 或 \sbl
顺时针循环 $W $\sbl $0$ ；逆时针循环 $W $\sbl $0$

热机效率：$\eta =$\sbl $\times 100\%$

热力学第二定律：
热量不能自发从\sbl 物体传到\sbl 物体
不可能从单一热源吸收热量，使其全部用来做功，而不引起其它影响

气缸问题解题步骤：
1. 确定\sbl 态
2. 写 \sbl 关系式
3. 写 $p-V-T$ 关系式
4. 求解

求压强 $P$：对自由活动的部分进行\sbl 分析，求得大气压强，得压强
mmHg 是\sbl 单位

\subsection*{国际单位制 (SI)}

七个基本单位：
米 (\sbl )
千克 (\sbl )
秒 (\sbl )
安培 (\sbl )
开尔文 (\sbl )
摩尔 (\sbl )
坎德拉 (\sbl )

注意：电压 (V) 和力 (F) \sbl 基本单位
F 也可以是\sbl 的单位

\subsection*{电路补充}

电动势定义：非静电力把 \sbl C 正电荷从负极搬到正极所做的功

电源输出功率最大条件：$R_{\text{外}} =$\sbl $r_{\text{内}}$

电动机能量转化：存在 $I^2R$ 的\sbl 损

串并联特点：
串联：$I$ \sbl ，$U \propto$ \sbl
并联：$V$ \sbl ，$I \propto$ \sbl

电源内阻影响：
改变滑动变阻器，与 $R$ 并联部分的 $U/I$ 变化趋势\sbl （有内阻时）
电源无内阻时，仅限与变阻器同一个支路上 $U/I$ 变化趋势\sbl

并联电阻极值：两个支路阻值尽可能\sbl 时，总电阻最大

欧姆定律相关：
由 $U/I \Rightarrow$ \sbl 定律
$\Delta U/\Delta I \Rightarrow$ 伏安特性曲线斜率 $=$ \sbl $r$

上一个电路 $R$ 最大时 $R_{\text{总}}$ \sbl

\subsection*{力学细节补充}

绳结模型：
活绳弹力处处\sbl ；死结绳弹力\sbl
活杆支持力沿\sbl ；固定杆支持力沿\sbl 方向

自由落体到弹簧上：
当 $F_{\text{弹}} = mg$ 时，速度 $v $\sbl $v_{\max}$ ，加速度 $a =$\sbl
最低点：加速度 $a_{\max}$ ，$a $\sbl $g$ ，速度 $v =$\sbl

弹簧连接体突变问题：
弹簧吊着小球，小球用轻绳吊着另一个小球，剪断轻绳后：
弹簧长度\sbl 突变，弹力\sbl 突变